
\section{单元测试详细说明}

\subsection{ECS 核心验证}

\texttt{ecsCorePhase1.verify.ts} 覆盖：创建多个实体并验证 ID 唯一；添加与移除组件后 \texttt{getComponent} 行为；销毁实体后组件不可访问；注册两个不同优先级 System 时执行顺序符合分组与 priority。该脚本可在 CI 中作为门禁，防止内核回归\cite{petty2010vv}。

\subsection{公式解析器验证}

\texttt{formulaParser.verify.ts} 覆盖：四则运算、括号、属性别名替换；除零与非法字符应抛出或返回安全默认值（以实现为准）；空公式与超长公式边界。技能伤害若依赖公式，任何解析器变更须 rerun 此脚本。

\subsection{属性修饰器验证}

\texttt{attributeModifier.verify.ts} 覆盖：同属性多个 modifier 叠加顺序；按 \texttt{sourceId} 批量移除后恢复 base；\texttt{getModifierContributions} 返回可追溯列表。升级奖励与技能 Buff 均依赖该机制，回归意义重大\cite{briand2001maintainability}。

\section{集成测试场景扩展}

\subsection{跑图可达性测试}

对 \texttt{RunMapGenerator.generate} 运行 100 次随机种子，断言每次 \texttt{validateActReachBoss} 为真；记录失败次数，应接近 0（仅 fallback 路径例外）。可将种子与 Act 索引打印，便于复现失败样例\cite{gregory2018game}。

\subsection{装配同步测试}

步骤：进入战斗 $\rightarrow$ Tab 打开面板 $\rightarrow$ 将技能 A 与 B 互换 $\rightarrow$ 关闭面板 $\rightarrow$ 观察施法图标与伤害类型是否随 B 变化。预期：下一冷却周期使用新技能。若立即仍用 A，则 \texttt{SkillLoadoutSyncSystem} 存在 bug。

\subsection{Worker 一致性测试}

在相同随机种子与输入序列下，分别运行 Worker 开/关各 30 秒，对比怪物位置轨迹是否一致（允许一帧延迟）。若偏差过大，说明主线程与 Worker 逻辑漂移，须统一 \texttt{combatWorkerLogic.ts}\cite{mckenzie2008crowd}。

\section{测试数据与截图要求}

终稿建议包含不少于 6 张截图：主菜单、跑图界面、战斗场面、技能面板、升级奖励、重启面板。每张图配 100--200 字说明。性能测试附 1 张 Performance 时间轴截图。缺少截图会降低论文说服力。

\section{测试结论补充说明}

综合第~\ref{chap:test}~章主文与上文扩展用例，可得出如下结论：系统在功能上满足 Must 级需求；Should 级需求中跑图 DAG、技能装配 UI、升级奖励与对象池均已通过手工用例；Could 级法杖三槽完整 UI 与部分 Effect 独立特效尚未实现，与第~\ref{chap:conclusion}~章展望一致。性能方面，在 Windows + Laya IDE 内置 Chromium、1920$\times$1080 预览环境下，Level3 五波实验整场平均 FPS 约 34.48；千怪同窗消融表明对象池将 P95 帧时间由 76.50\,ms 降至 48.60\,ms（约降 36.5\%），而平均 FPS 提升有限。动态图集在千怪波次上带来约 13.4\% 的 FPS 提升。单元验证脚本 \texttt{ecsCorePhase1.verify.ts}、\texttt{formulaParser.verify.ts}、\texttt{attributeModifier.verify.ts} 在提交前均通过。遗留缺陷以配置同步、部分特效占位为主，不影响答辩演示主路径。
